\subsection{Change of Rings}

Let A and B be rings, and let \(f : A \to B\) be a ring homomorphism. Let C be an additive set of A-modules and D an additive set of B-modules.

First, suppose that for every B-module M of type D, the A-module \(f_{*}(M)\) obtained by restricting the ring of scalars to A is of type C. Then the mapping from D to \(\mathrm{K}(\mathcal{C})\) that sends M to \([f_{*}(M)]_{\mathcal{C}}\) is an additive function of modules; from it, we deduce a group homomorphism

\[
f _ {*}: \mathrm{K} (\mathcal {D}) \longrightarrow \mathrm{K} (\mathcal {C}).
\]

We define a group homomorphism \(f_{*}:\mathrm{K}^{\prime}(\mathcal{D})\to \mathrm{K}^{\prime}(\mathcal{C})\) likewise.

We now suppose that for every A-module E of type C, the B-module \(f^{*}(\mathrm{E})\) derived from E by extension of scalars via f (II, §5, No.~1, p.~277) is of type D. The mapping from C to \(\mathrm{K}'(\mathcal{D})\) that sends an element E of C to \([f^{*}(\mathrm{E})]_{\mathcal{D}}'\) is a weakly additive function of modules; it therefore induces a group homomorphism \(f^{*}: \mathrm{K}'(\mathcal{C}) \to \mathrm{K}'(\mathcal{D})\) .

Suppose, furthermore, that for every exact sequence

\[
0 \longrightarrow \mathrm{E} ^ {\prime} \longrightarrow \mathrm{E} \longrightarrow \mathrm{E} ^ {\prime \prime} \longrightarrow 0
\]

of A-modules of type C, the sequence of B-modules of type D

\[
0 \longrightarrow \mathrm{B} \otimes_ {\mathrm{A}} \mathrm{E} ^ {\prime} \longrightarrow \mathrm{B} \otimes_ {\mathrm{A}} \mathrm{E} \longrightarrow \mathrm{B} \otimes_ {\mathrm{A}} \mathrm{E} ^ {\prime \prime} \longrightarrow 0
\]

is exact. This holds, in particular, in the following cases:

\begin{enumerate}
\def\labelenumi{\alph{enumi})}
\item
  The homomorphism \(f\) makes B into a projective A-module (II, §3, No.~6, p.~251, Proposition 5 and II, §3, No.~7, p.~257, Corollary 6) *or, more generally, a flat (X, §1, n° 3, p.~8, définition 1) A-module*.
\item
  *The set C is a set of classes of projective or, more generally, flat (X, §4, n° 5, p.~72, corollaire 2) A-modules*.
\end{enumerate}

The mapping from \(\mathcal{C}\) to \(\mathrm{K}(\mathcal{D})\) that sends \(\operatorname{E}\) to \([f^{*}(\operatorname{E})]_{\mathcal{D}}\) is then additive. It therefore induces a group homomorphism \(f^{*}:\mathrm{K}(\mathcal{C})\to \mathrm{K}(\mathcal{D})\) .
% semantic-fragments: legacy-input-seam
\relax{}
